\section{Implementação em R}
\label{sec:medidas em r}

As medidas vistas neste capítulo estão, em grande parte, disponíveis nas funções básicas
da linguagem R, de modo que seu cálculo, na prática, se resume a poucas linhas de código.
Ainda assim, é preciso conhecer as convenções adotadas por cada função, pois nem sempre
elas coincidem com as definições apresentadas aqui -- como é o caso do divisor da
variância e da constante do desvio absoluto mediano.
Nesta seção, veremos como computar cada uma delas, e, quando a linguagem não oferecer uma
função pronta, como implementá-la a partir da definição.

Usaremos, ao longo dos exemplos, o vetor
\[
    x = (0, 10, 10, 10, 11, 11, 12, 13, 14),
\]
o mesmo da seção de simetria, de sorte que os valores obtidos possam ser conferidos com
aqueles já calculados.
As saídas de cada comando são mostradas como comentários logo abaixo dele.

\subsection{Medidas de Posição}

O R não possui uma função para a moda amostral -- a função \texttt{mode} retorna o modo de um objeto,
como \texttt{"numeric"} ou \texttt{"character"}, e não a medida estatística.
Pela definição de moda amostral, basta tomar o valor que maximiza a frequência
absoluta $N(x)$, o que pode ser feito a partir da tabela de frequências produzida por
\texttt{table}.
Como o máximo pode não ser único, o código \ref{cod:moda} retorna todos os valores que o
atingem, o que deixa explícito quando a distribuição é multimodal.

\begin{lstlisting}[language=R, caption={Moda Amostral}, label={cod:moda}]
moda <- function(x) {
    freq <- table(x)
    as.numeric(names(freq)[freq == max(freq)])
}

x <- c(0, 10, 10, 10, 11, 11, 12, 13, 14)
moda(x)
# [1] 10
\end{lstlisting}

Note que \texttt{table} converte os valores em nomes, isto é, em cadeias de caracteres,
daí a necessidade de \texttt{as.numeric}; para atributos qualitativos, essa conversão deve
ser omitida.

As estatísticas de ordem são obtidas com \texttt{sort}, e a mediana e a média, com
\texttt{median} e \texttt{mean} respectivamente.
A função \texttt{median} segue a definição \ref{def:mediana amostral}, tomando a média
das duas estatísticas centrais quando $n$ é par.

\begin{lstlisting}[language=R, caption={Estatísticas de Ordem, Mediana e Média}, label={cod:posicao}]
sort(x)
# [1]  0 10 10 10 11 11 12 13 14
median(x)
# [1] 11
mean(x)
# [1] 10.11111
\end{lstlisting}

Se houver algum valor ausente no vetor, \texttt{median} e \texttt{mean} retornam
\texttt{NA}, assim como \texttt{var}, \texttt{sd} e \texttt{mad}, vistas adiante; já
\texttt{quantile} e \texttt{IQR} interrompem a execução com um erro.
Para descartar os valores ausentes antes do cálculo, usa-se, em todas essas funções, o
argumento \texttt{na.rm = TRUE}.
Por outro lado, \texttt{sort} e \texttt{table} -- e, portanto, a função \texttt{moda} --
descartam os valores ausentes sem nenhum aviso, o que pode esconder o fato de que o valor
ausente é o mais frequente do atributo.

\subsection{Medidas de Dispersão}

A variância e o desvio-padrão são dados pelas funções \texttt{var} e \texttt{sd}, ambas
com divisor $n - 1$, conforme a definição de variância amostral.
Caso se deseje a variância sob a probabilidade empírica, isto é, com divisor $n$, basta
multiplicar o resultado por $\frac{n-1}{n}$.
Já o desvio absoluto médio e o coeficiente de variação não possuem funções próprias, mas
são escritos diretamente a partir de suas fórmulas.

\begin{lstlisting}[language=R, caption={Medidas de Dispersão}, label={cod:dispersao}]
n <- length(x)

mean(abs(x - mean(x)))    # desvio absoluto medio
# [1] 2.320988
var(x)
# [1] 16.36111
var(x) * (n - 1) / n      # variancia com divisor n
# [1] 14.54321
sd(x)
# [1] 4.044887
sd(x) / mean(x)           # coeficiente de variacao
# [1] 0.4000438
\end{lstlisting}

\subsection{Quantis}

Os quantis empíricos são computados pela função \texttt{quantile}, que recebe o vetor de
observações e um vetor \texttt{probs} de probabilidades.
Quando este é omitido, são retornados os quartis junto do mínimo e do máximo.
A função implementa nove métodos distintos de determinação de quantis, escolhidos pelo
argumento \texttt{type}; o padrão, \texttt{type = 7}, é justamente o método de
interpolação do algoritmo \ref{alg:quantil}.
O intervalo interquartílico, por sua vez, é dado por \texttt{IQR}, que usa o mesmo método.

\begin{lstlisting}[language=R, caption={Quantis e Intervalo Interquartílico}, label={cod:quantis}]
quantile(x)
#   0%  25%  50%  75% 100%
#    0   10   11   12   14
quantile(x, probs = c(0.05, 0.95))
#   5%  95%
#  4.0 13.6
IQR(x)
# [1] 2
\end{lstlisting}

Para conferir a equivalência, o código \ref{cod:quantil interpolacao} traduz o algoritmo
\ref{alg:quantil} para o R.
Como os vetores do R são indexados a partir de $1$, a tradução é direta.

\begin{lstlisting}[language=R, caption={Quantil Empírico por Interpolação}, label={cod:quantil interpolacao}]
quantil <- function(x, p) {
    s <- sort(x)
    n <- length(s)
    if (n == 1) return(s[1])
    h <- (n - 1) * p + 1
    j <- min(floor(h), n - 1)
    gama <- h - j
    (1 - gama) * s[j] + gama * s[j + 1]
}

sapply(c(0.05, 0.25, 0.5, 0.75, 0.95), quantil, x = x)
# [1]  4.0 10.0 11.0 12.0 13.6
\end{lstlisting}

A função \texttt{summary}, aplicada a um vetor numérico, reúne num só comando o mínimo,
os quartis, a média e o máximo, sendo um bom primeiro resumo de um atributo quantitativo.
Deve-se ter cuidado, entretanto, com a função \texttt{fivenum}, que retorna as dobradiças
de Tukey em vez dos quartis do método padrão; as duas coincidem no exemplo, mas diferem
em geral.

\begin{lstlisting}[language=R, caption={Resumo com \texttt{summary} e \texttt{fivenum}}, label={cod:summary}]
summary(x)
#    Min. 1st Qu.  Median    Mean 3rd Qu.    Max.
#    0.00   10.00   11.00   10.11   12.00   14.00
fivenum(x)
# [1]  0 10 11 12 14
\end{lstlisting}

O \textit{box-plot} é construído pela função \texttt{boxplot}, e os valores usados no seu
desenho podem ser obtidos com \texttt{boxplot.stats}.
Seu componente \texttt{out} contém as observações fora dos bigodes, isto é, os possíveis
valores atípicos.
Os quartis usados por essas funções também são as dobradiças de Tukey.

\begin{lstlisting}[language=R, caption={Box-plot e Valores Atípicos}, label={cod:boxplot}]
boxplot(x, horizontal = TRUE, xlab = "x")

boxplot.stats(x)$out
# [1] 0
\end{lstlisting}

No exemplo, $q_1 - 1.5\, d_q = 10 - 3 = 7$, e a observação $0$ é a única abaixo desse
limite.

\subsection{Simetria e Assimetria}

Nenhum dos coeficientes de assimetria vistos está presente nas funções básicas do R.
O código \ref{cod:assimetria} os implementa a partir das fórmulas da seção de simetria,
reaproveitando a função \texttt{moda} do código \ref{cod:moda}.
O coeficiente de Pearson I usa apenas a primeira moda retornada, pois só faz sentido
quando a distribuição é unimodal.

\begin{lstlisting}[language=R, caption={Coeficientes de Assimetria}, label={cod:assimetria}]
pearson_1 <- function(x) (mean(x) - moda(x)[1]) / sd(x)

pearson_2 <- function(x) 3 * (mean(x) - median(x)) / sd(x)

quartilico <- function(x) {
    q <- quantile(x, c(0.25, 0.5, 0.75), names = FALSE)
    (q[1] + q[3] - 2 * q[2]) / (q[3] - q[1])
}

fisher <- function(x) {
    n <- length(x)
    sigma <- sqrt(mean((x - mean(x))^2))
    sqrt(n * (n - 1)) / (n - 2) * sum((x - mean(x))^3) / (n * sigma^3)
}

c(pearson_1(x), pearson_2(x), quartilico(x), fisher(x))
# [1]  0.02746952 -0.65926851  0.00000000 -2.28919119
\end{lstlisting}

Os valores reproduzem o exemplo em que os coeficientes discordam em sinal.
Note que, em \texttt{fisher}, o desvio-padrão $\sigma$ é calculado com divisor $n$, e
não com \texttt{sd}.
O mesmo coeficiente é oferecido pela função \texttt{skewness} do pacote \texttt{e1071},
com o argumento \texttt{type = 2}.

Por fim, o gráfico de simetria é construído tomando os pares $(u_i, v_i)$ a partir das
estatísticas de ordem e traçando a reta $v = u$, como no código \ref{cod:grafico simetria}.

\begin{lstlisting}[language=R, caption={Gráfico de Simetria}, label={cod:grafico simetria}]
grafico_simetria <- function(x) {
    s <- sort(x)
    n <- length(s)
    m <- median(s)
    i <- seq_len(n %/% 2)
    u <- m - s[i]
    v <- s[n - i + 1] - m
    plot(u, v, pch = 19, xlab = "u", ylab = "v",
         xlim = c(0, max(u, v)), ylim = c(0, max(u, v)))
    abline(a = 0, b = 1, lty = 2)
}

grafico_simetria(x)
\end{lstlisting}

Para o vetor $x$, os pares obtidos são $(11, 3)$, $(1, 2)$, $(1, 1)$ e $(1, 0)$.
O primeiro deles, que retrata as caudas, fica bem abaixo da reta, o que evidencia a
cauda inferior longa causada pela observação $0$.

\subsection{Medidas Robustas}

A média aparada é obtida pela própria função \texttt{mean}, por meio do argumento
\texttt{trim}, que corresponde ao parâmetro $\alpha$ da definição de média aparada:
são descartadas as $\lfloor n\alpha \rfloor$ menores e maiores observações.
O desvio absoluto mediano é dado pela função \texttt{mad}, mas, por padrão, ela multiplica
o resultado pela constante $1.4826$, que o torna comparável ao desvio-padrão quando os dados
provêm de uma distribuição normal.
Para obter a medida tal como definida aqui, deve-se usar \texttt{constant = 1}.

Retomemos o exemplo dos salários com um vetor de dez salários, em milhares de reais, no
qual há um único valor muito acima dos demais.

\begin{lstlisting}[language=R, caption={Medidas Robustas}, label={cod:robustas}]
salario <- c(2.1, 2.3, 2.4, 2.6, 2.8, 3.0, 3.1, 3.5, 4.0, 25.0)

c(mean(salario), median(salario), mean(salario, trim = 0.1))
# [1] 5.0800 2.9000 2.9625
c(sd(salario), IQR(salario), mad(salario, constant = 1))
# [1] 7.02263 0.95000 0.55000
mad(salario)
# [1] 0.81543
\end{lstlisting}

A média é quase o dobro da mediana, ao passo que a média aparada com $\alpha = 0.1$, que
descarta apenas a menor e a maior observação, já se aproxima desta.
Da mesma forma, o desvio-padrão é muito maior que o intervalo interquartílico e o desvio
absoluto mediano, pois é dominado pelo valor $25$.

\subsection{Resumo por Grupos}

Na prática, as medidas costumam ser computadas para cada nível de um atributo
qualitativo do quadro, a fim de comparar as distribuições entre os grupos.
O pacote \texttt{dplyr} permite fazê-lo agrupando o quadro com \texttt{group\_by} e
aplicando as funções vistas dentro de \texttt{summarise}, como no código
\ref{cod:resumo grupos}.
O resultado de \texttt{summarise} é um \textit{tibble}, que arredonda os valores ao ser
impresso; a conversão com \texttt{as.data.frame} serve apenas para exibi-los com mais
casas decimais.

\begin{lstlisting}[language=R, caption={Medidas de Resumo por Grupos}, label={cod:resumo grupos}]
library(dplyr)

salarios <- data.frame(
    setor = rep(c("A", "B"), each = 6),
    salario = c(2.1, 2.4, 2.6, 2.8, 3.0, 3.1,
                2.3, 2.5, 3.5, 4.0, 4.2, 25.0)
)

salarios |>
    group_by(setor) |>
    summarise(media = mean(salario),
              mediana = median(salario),
              aparada = mean(salario, trim = 0.2),
              dp = sd(salario),
              iqr = IQR(salario),
              mad = mad(salario, constant = 1)) |>
    as.data.frame()
#   setor    media mediana aparada        dp iqr  mad
# 1     A 2.666667    2.70    2.70 0.3777124 0.5 0.30
# 2     B 6.916667    3.75    3.55 8.8925624 1.4 0.85
\end{lstlisting}

No setor A, as medidas clássicas e as robustas praticamente coincidem, enquanto, no setor
B, a média e o desvio-padrão são inflados por um único salário, o que as medidas robustas
revelam de imediato.
